\section{Related Work}
\label{sec:related}

\subsection{Traditional Handwriting Verification}
Early writer verification relied on handcrafted descriptors extracted from
the handwritten trace. Three families dominated: \emph{grapheme-based}
methods, which model the statistical distribution of elementary stroke
shapes~\cite{bensefia2005}; \emph{texture-based} methods, which treat
handwriting as an oriented texture and encode it with Gabor filters or
grey-level co-occurrence statistics~\cite{said2000}; and
\emph{contour-based} methods, which measure slant, curvature and shape
along ink boundaries~\cite{schomaker2004}. Textural and
allographic features were combined into strong text-independent baselines by
Bulacu and Schomaker~\cite{bulacu2007}, within the broader framing of the
handwriting-recognition problem surveyed by Plamondon and
Srihari~\cite{plamondon2000}. These approaches perform well in controlled
settings but degrade under writing variabilityand noise and
require expert feature engineering that transfers poorly across datasets.

\subsection{Deep Learning for Writer Analysis}
Convolutional Neural Networks (CNNs) replaced hand-designed descriptors with
hierarchical features learned directly from pixels. In writer identification,
Fiel and Sablatnig~\cite{fiel2015} used CNN activations as a per-writer
feature vector, while Christlein et al.~\cite{christlein2015} employed
penultimate-layer activations as local descriptors aggregated into a global
GMM supervector; the same group subsequently removed the need for writer
labels altogether, learning descriptors by unsupervised training on
surrogate classes~\cite{christlein2017}. Later work moved to deeper,
fragment-level representations, as in FragNet~\cite{he2020fragnet}. These
systems, however, target \emph{closed-set identification}: they assign a
probe to one of a fixed set of enrolled writers and are evaluated with
top-$k$ accuracy or mean average precision (mAP).

They do not
address the \emph{open-set verification} question—whether two arbitrary
samples share an author—which requires an explicit similarity model and a
decision threshold rather than a closed classifier, and which is the setting
targeted here.

\subsection{Siamese Networks and Metric Learning}
Verification is naturally cast as metric learning: a Siamese architecture
applies two weight-sharing encoders to an input pair and learns an embedding
in which same-writer samples are close and different-writer samples are far
apart. The paradigm originates with Bromley et al.~\cite{bromley1994siamese},
who introduced twin networks for signature verification, and was later
popularised for one-shot recognition by Koch et al.~\cite{koch2015siamese}.
Two metric objectives shape the embedding geometry explicitly: the Contrastive
Loss of Hadsell et al.~\cite{hadsell2006contrastive}, which pulls genuine
pairs together and pushes impostor pairs beyond a margin, and the Triplet Loss
of Schroff et al. ~\cite{schroff2015facenet}, which enforces a \emph{relative}
constraint—an anchor must be closer to a positive than to a negative by a
margin—yielding embeddings well suited to open-set biometrics. In handwriting,
triplet CNNs were applied to writer retrieval by
Keglevic et al.~\cite{keglevic2018}. A third, non-metric option keeps the
Siamese encoder but trains a decision head directly on the pair with Binary
Cross-Entropy (BCE), as in the one-shot formulation of Koch et
al.~\cite{koch2015siamese}; this remains the conventional
pairwise-classification baseline. The three objectives therefore compete for
the same architecture, each structuring the embedding space differently: BCE
optimises a decision head without explicitly constraining the space, whereas
the two metric losses act on pairwise geometry directly. To the best of our
knowledge, no prior study compares these three objectives under a single
fixed training recipe across a shared set of modern backbones on handwriting
data, and that is the gap this work addresses.

\subsection{Backbone Architectures and Transfer Learning}
The encoder choice governs the quality of the learned representation. Residual
connections~\cite{he2016resnet} enabled very deep networks;
EfficientNet~\cite{tan2019efficientnet} reached comparable accuracy at a
fraction of the parameters through compound scaling; and
MobileNetV3~\cite{howard2019mobilenetv3} prioritised efficiency via
depthwise-separable convolutions. Fine-tuning
ImageNet-pretrained~\cite{deng2009imagenet} backbones is now standard, but the
domain gap between natural images and near-binary handwriting makes the
optimal backbone and transfer strategy for verification non-obvious—motivating
the controlled backbone sweep reported here.

\subsection{Evaluation of Biometric Verification}
As a biometric authentication task, writer verification must be evaluated with
detection-oriented metrics rather than plain accuracy. The Equal Error Rate
(EER), the operating point at which the False Acceptance Rate equals the False
Rejection Rate, provides a single-number, threshold-independent summary, while
the ROC curve and its area (AUC) characterise performance across all
thresholds; decidability ($d'$) quantifies the separation between the genuine
and impostor score distributions~\cite{mansfield2002}. We adopt this protocol
throughout and additionally report operation at fixed low-FAR points, the
regime relevant to deployed systems and not captured by the EER alone.